\section*{Chapitre \ref{ch:in:portfolio-manager-and-pod-analyst} --- Gérant de portefeuille et analyste de pod}

\begin{solution}{exo:in:portfolio-manager-and-pod-analyst:1}
La perte reportée est d'abord absorbée~: $40-10=30$ millions de profit net, et 15~\% de ce montant représentent 4.5~millions de dollars, dont le salaire de 0.5
million de dollars déjà versé est une avance. Total pour l'année~: 4.5~millions de dollars.
\end{solution}

\begin{solution}{exo:in:portfolio-manager-and-pod-analyst:2}
La rémunération espérée augmente de 7.30~millions de dollars (de 3.18 à 10.48~millions de dollars)~; la probabilité d'arrêt baisse de 42.1~points (de 57.3~\% à
15.2~\%).
\end{solution}

\begin{solution}{exo:in:portfolio-manager-and-pod-analyst:3}
Environ $2\sqrt3=3.46$~; avec la correction de Lo pour 756~observations quotidiennes, 3.45.
\end{solution}

\begin{solution}{exo:in:portfolio-manager-and-pod-analyst:4}
Le versement est une option d'achat sur le profit de chaque année~: le gérant partage les bonnes années et ne perd rien les mauvaises, mais les versements
liés aux pertes reportées se compensent, si bien qu'un pod d'espérance nulle a tout de même un versement espéré positif. Les investisseurs le paient, comme
le coût de compensation du livre~16, chapitre~3~; l'arrêt le limite en mettant fin aux ancienneté malheureuses.
\end{solution}

\begin{solution}{exo:in:portfolio-manager-and-pod-analyst:5}
L'équivalent certain pénalise la dispersion des issues~: un gérant averse au risque valorise davantage les nombreuses issues faibles que
les rares issues élevées, il lui faut donc plus de compétence avant que l'accord vaille un salaire certain.
\end{solution}

\begin{solution}{exo:in:portfolio-manager-and-pod-analyst:6}
La volatilité du sous-portefeuille et son ratio de Sharpe espéré (qui fixent la distribution du versement~: à 6~\% de volatilité,
3~millions de dollars de profit annuel par unité de ratio de Sharpe, 300\,000~dollars à 10~\%)~; les limites et la question de savoir si un arrêt met fin à son poste~;
si les pertes sont reportées~; ce qu'a été le bonus pour les analystes de son niveau~; et ce que vaut son propre historique
pour son prochain poste.
\end{solution}

\begin{solution}{exo:in:portfolio-manager-and-pod-analyst:7}
Pour $S=0$, la probabilité d'arrêt passe de 57.3~\% à 26.8~\% et la rémunération espérée passe de 3.18 à 3.56~millions de dollars~; pour
$S=1$, la probabilité d'arrêt passe de 15.2~\% à 2.1~\% et la rémunération passe de 10.48 à 10.76~millions de dollars. Après la réduction à
5~\%, le capital est divisé par deux, si bien qu'une nouvelle perte de 5~\% du capital initial représente une perte de 10~\% sur ce qui reste.
\end{solution}

\begin{solution}{exo:in:portfolio-manager-and-pod-analyst:8}
Un ratio de Sharpe de 1.2 sur trois ans a une statistique de Student d'environ $1.2\sqrt3=2.08$, à peine significative~; elle faisait partie de
nombreux gérants, et ceux qui ont été arrêtés ne figurent pas dans l'échantillon~; 43~\% des ancienneté sans compétence survivent à trois ans dans
l'exemple du chapitre. Doubler le capital modifie aussi les coûts de la stratégie et sa capacité. Il faut examiner le processus, le risque pris
et le degré de concentration du résultat, puis augmenter le capital par paliers.
\end{solution}

\begin{solution}{pb:in:portfolio-manager-and-pod-analyst:1}
\begin{enumerate}
\item Comme dans la définition du chapitre.
\item Un pourcentage du profit annuel propre au gérant, valorisé au prix de marché, sans égard aux autres gérants~; pertes
reportées~; salaire considéré comme une avance~; pas de récupération~; versement dû même si les fonds perdent.
\item Il alloue et réalloue le capital entre les gérants~; le capital fixe la taille du pod, donc son profit et son risque.
\item Une réduction du capital au premier drawdown et un arrêt au second~; livre~16, chapitre~8.
\item Les coûts propres du pod, y compris la rémunération de son équipe et les données.
\item $\max(w,\rho\max(P_t-L_{t-1},0)K)$ chaque année avec $L_t=\max(L_{t-1}-P_t,0)$.
\item Convexe en fonction du profit (une option d'achat)~; interrompue par l'arrêt (une option à barrière désactivante).
\item 500~millions de dollars, volatilité de 6~\%, versement de 15~\%, salaire de 0.5~million de dollars, réduction à 5~\%, arrêt à 7.5~\%, trois ans.
\item 3.18~millions de dollars et 57.3~\%~; 10.48~millions de dollars et 15.2~\%~; 23.36~millions de dollars et 1.9~\%.
\item La somme certaine ayant la même utilité espérée~; aversion relative au risque égale à 3 et 5~millions de dollars d'autre patrimoine.
\item 3.85~ans de résultats quotidiens.
\item La capacité, l'encombrement (crowding), le risque pris, la concentration du résultat, le processus et la question de savoir si l'équipe suit le gérant.
\item Un mandat face à un indice de référence, évalué par rapport à celui-ci, sans arrêt dans l'accord.
\item Salaires médians de 223\,860~dollars pour les gérants financiers et de 124\,370~dollars pour les analystes du secteur des valeurs mobilières~; ils laissent
de côté les versements et les bonus.
\item Par l'intermédiaire d'un sous-portefeuille et d'un historique qui lui est propre.
\item $S\approx0.26$, avec environ 43~\% des ancienneté arrêtées.
\item Avec une aversion au risque de 3 et 5~millions de dollars de patrimoine, le point de bascule se déplace vers $S\approx0.75$ (environ 23~\% d'arrêts).
\item De 6.04~millions de dollars et 32.5~\% d'arrêts pour $S=0.5$ à 10.48~millions de dollars et 15.2~\% pour $S=1$.
\item Une garantie pour la première année, la taille de l'arrêt, la question de savoir si le capital est réduit avant l'arrêt, et ce qu'il advient des pertes
reportées si elle quitte le poste.
\item Si son ratio de Sharpe est réellement compris entre 0.5 et 1, la rémunération espérée de l'accord représente deux à trois fois son salaire, même corrigée de l'aversion
au risque, mais dans un tiers à un sixième des ancienneté elle sera hors du poste dans les trois ans. Il ne faut l'accepter qu'avec une épargne suffisante pour absorber
ce scénario et une garantie pour la première année.
\end{enumerate}
\end{solution}

\begin{solution}{iq:in:portfolio-manager-and-pod-analyst:1}
Réduire le risque avant que l'échelle de réduction ne s'en charge, conserver les positions offrant le meilleur rendement espéré par unité de risque, vérifier
si la perte vient du processus ou d'une seule position, et informer le responsable des risques de ce que l'on fait.
\iqlookfor{réduction délibérée du risque et communication, non l'espoir.}
\end{solution}

\begin{solution}{iq:in:portfolio-manager-and-pod-analyst:2}
Avec les éléments de preuve~: depuis quand cela a cessé de fonctionner, comment on le sait, ce que cela a coûté et ce qu'on propose à la place, le plus tôt possible.
\iqlookfor{franchise appuyée sur des preuves, et tôt.}
\end{solution}

\begin{solution}{iq:in:portfolio-manager-and-pod-analyst:3}
Avec des pertes reportées, le versement est très hors de la monnaie, et davantage de risque en accroît la valeur~: le détenteur d'une option apprécie
la volatilité. Les plateformes répondent par des limites de risque, des réductions de capital et des arrêts, qui plafonnent le risque qu'un gérant peut ajouter.
\iqlookfor{convexité de l'option et limites qui la compensent.}
\end{solution}

\begin{solution}{iq:in:portfolio-manager-and-pod-analyst:4}
Régresser les rendements quotidiens du pod sur les rendements du marché et des facteurs~; les expositions multipliées par les rendements des facteurs constituent la part de bêta.
Les plateformes souhaitent souvent qu'elle soit proche de zéro et peuvent la couvrir centralement.
\iqlookfor{régression factorielle et raison pour laquelle la plateforme s'en soucie.}
\end{solution}

\begin{solution}{iq:in:portfolio-manager-and-pod-analyst:5}
Les rendements quotidiens, le risque et les limites dans lesquels ils ont été obtenus, la capacité, la concentration, le taux de rotation, l'indice de référence et
les expositions aux facteurs, et la question de savoir si l'équipe et les données suivent le gérant.
\iqlookfor{les statistiques de l'historique et ses conditions.}
\end{solution}

\begin{solution}{iq:in:portfolio-manager-and-pod-analyst:6}
$0.15\,E[\max(X,0)]=0.15\,\sigma/\sqrt{2\pi}=0.15\times30/2.507=1.80$~millions de dollars.
\iqlookfor{l'espérance de la partie positive d'une loi normale.}
\end{solution}
